\documentclass[fontsize=12pt]{mybib}

\setheaderlogo{logo.png}

\title{Der zweite Tempel}
\subtitle{Periode von 515\vChr{} bis 70\nChr}
\author{theologicon.de}
\date{2026}

\renewcommand{\DFA}{15mm}
\renewcommand{\DZA}{1.0}
\setlength{\footheight}{30.4pt}

\begin{document}
\maketitle
\tableofcontents

\section{Einleitung}

Das Judentum zur Zeit des Zweiten Tempels umfasst die Epoche zwischen der Wiedererrichtung des Jerusalemer Tempels um 515/516 v. Chr. und seiner Zerstörung durch die Römer im Jahr 70 n. Chr. In diesen Jahrhunderten veränderten sich die religiösen, gesellschaftlichen und politischen Lebensbedingungen der jüdischen Bevölkerung unter persischer, hellenistischer und römischer Herrschaft. Zugleich entwickelten sich unterschiedliche Formen jüdischer Schriftauslegung, religiöser Praxis und theologischer Reflexion. Diese Entwicklungen bildeten einen wesentlichen historischen Hintergrund für die spätere Entstehung des rabbinischen Judentums und des frühen Christentums (Sanders 1992; Grabbe 2010).

Die Bezeichnung „Zweites Tempel-Judentum“ beschreibt dabei mehr als einen durch zwei historische Ereignisse begrenzten Zeitraum. Sie verweist auf eine Phase religiöser Neuordnung, deren Voraussetzungen bereits im babylonischen Exil lagen. Die Zerstörung des Ersten Tempels und der Verlust des judäischen Königtums hatten die bisherige Verbindung von Land, politischer Herrschaft und Tempelkult verändert. Mit der Rückkehr eines Teils der Exilierten, dem Wiederaufbau des Tempels und der Neuorganisation der Gemeinschaft wurden ältere Traditionen aufgenommen und unter veränderten Bedingungen weiterentwickelt. Dabei entstanden Formen religiöser Zugehörigkeit, die nicht ausschließlich an das Land Israel oder die unmittelbare Teilnahme am Jerusalemer Tempelkult gebunden waren (Grabbe 2010).

Der Jerusalemer Tempel blieb bis zu seiner Zerstörung ein religiöser und gesellschaftlicher Mittelpunkt. Daneben prägten die Tora und ihre Auslegung, das gemeinschaftliche Gebet, die religiöse Unterweisung und das Leben in der Diaspora die jüdische Identität. Innerhalb dieser gemeinsamen Traditionswelt bestanden unterschiedliche religiöse Strömungen, soziale Gruppen und theologische Auffassungen. Fragen der Gesetzesauslegung, der kultischen Reinheit, der göttlichen Offenbarung, der Auferstehung und der messianischen Hoffnung wurden keineswegs einheitlich beantwortet. Einblicke in diese Vielfalt vermitteln unter anderem die Schriftrollen vom Toten Meer, die apokalyptische und weisheitliche Literatur, die Werke Philos von Alexandria sowie die historischen Darstellungen des Flavius Josephus (Sanders 1992; Collins 2019).

Diese religiöse Vielfalt bildet zugleich den historischen Zusammenhang, in dem Jesus von Nazaret, seine ersten Anhänger und Paulus von Tarsus wirkten. Ihre Überzeugungen und Auseinandersetzungen entstanden innerhalb der jüdischen Welt ihrer Zeit. Das Judentum des Zweiten Tempels lässt sich daher nicht lediglich als Vorgeschichte des Christentums verstehen. Ebenso wenig war das spätere rabbinische Judentum bereits in der pharisäischen Bewegung vollständig vorgebildet. Beide Traditionen entwickelten sich über längere Zeiträume aus einer gemeinsamen, religiös vielfältigen Ausgangssituation und gewannen schrittweise eigene institutionelle und theologische Formen (Sanders 1992; Grabbe 2010; Collins 2019).

Der folgende Essay behandelt das Judentum der Zweiten Tempelzeit als eigenständigen Gegenstand der Religionsgeschichte. Untersucht werden seine Begriffs- und Forschungsgeschichte, die historischen Rahmenbedingungen, religiöse Gemeinschaften und Institutionen, die Entwicklung schriftlicher Überlieferungen sowie zentrale theologische Vorstellungen. Abschließend wird dargestellt, welche Bedeutung diese Epoche für die Entstehung des frühen Christentums und die Entwicklung des rabbinischen Judentums hatte.

\section{Begriffsklärung und Forschungsgeschichte}

Das Zweite Tempel-Judentum bezeichnet gewöhnlich die Zeit zwischen 515/516 v. Chr. und 70 n. Chr. (Grabbe 2010, 1–3; Collins 2019, 437–440). Der Begriff orientiert sich damit an einer zentralen religiösen Institution, die als sichtbares Zeichen der Gegenwart Gottes, als Opferstätte und als Mittelpunkt jüdischer Wallfahrt und Identität fungierte. Diese zeitliche Eingrenzung ist jedoch nicht mit einer scharfen religionsgeschichtlichen Grenze gleichzusetzen. Bereits während des babylonischen Exils hatten sich Formen jüdischen Lebens entwickelt, die ohne einen funktionierenden Jerusalemer Tempel auskommen mussten, und die theologischen und literarischen Traditionen des Zweiten Tempels endeten nicht mit dessen Zerstörung. Apokalyptische Schriften wie das vierte Esrabuch und die syrische Baruchapokalypse entstanden erst nach 70 n. Chr. und verarbeiteten gerade den Verlust des Heiligtums.

Die moderne Forschung verwendet deshalb auch den Begriff „Frühjudentum“, dessen zeitlicher Umfang je nach Fragestellung unterschiedlich bestimmt wird. Manche Darstellungen beginnen bereits mit der Exilszeit, andere konzentrieren sich auf die persische oder hellenistische Epoche und beziehen die Entwicklungen bis zum Bar-Kochba-Aufstand von 132–135 n. Chr. ein. Die Bezeichnung „Zweites Tempel-Judentum“ besitzt dennoch einen besonderen analytischen Wert, weil sie sichtbar macht, dass der Tempel die gesamte Epoche als religiöser Bezugspunkt prägte – selbst dort, wo jüdische Gemeinschaften weit entfernt von Jerusalem lebten oder den tatsächlichen Tempelkult kritisierten.

Die Forschungsgeschichte ist eng mit der Geschichte christlicher Bibelwissenschaft verbunden. Im 19. und frühen 20. Jahrhundert wurde häufig der Begriff „Spätjudentum“ verwendet, der in manchen wissenschaftlichen Zusammenhängen mit der Vorstellung verbunden war, das Judentum habe seine eigentliche religiöse Lebendigkeit verloren und sei zu einer erstarrten Gesetzesreligion geworden, die durch das Christentum überwunden worden sei (Schmid 2020, 139–152). Eine solche Deutung setzte das Christentum als Zielpunkt einer Entwicklung voraus und verkannte sowohl die religiöse Eigenständigkeit des Judentums als auch dessen innere Vielfalt. Die neuere Forschung hat diese Perspektive grundlegend korrigiert: Das Frühjudentum gilt heute nicht als Verfallsphase der alttestamentlichen Religion, sondern als eine Zeit intensiver theologischer Produktivität. Die Entstehung neuer Schriften, die Weiterentwicklung biblischer Traditionen, die Ausbildung unterschiedlicher religiöser Gemeinschaften und die Begegnung mit hellenistischer Philosophie bezeugen eine bemerkenswerte religiöse Dynamik.

Die Bezeichnung „Frühjudentum“ vermeidet die ältere Vorstellung eines geschichtlichen Niedergangs, ist aber ebenfalls nicht unproblematisch, weil sie den Eindruck erwecken kann, die jüdische Religion habe erst in dieser Phase begonnen. Zwischen dem vorexilischen Israel, der exilischen Gemeinschaft und dem späteren Judentum bestehen vielmehr historische Kontinuitäten und tiefgreifende Veränderungen zugleich. Für eine differenzierte Darstellung sind daher beide Begriffe sinnvoll: „Zweites Tempel-Judentum“ für die institutionell und chronologisch definierte Epoche und „Frühjudentum“ für den weiteren religionsgeschichtlichen Zusammenhang.

Ein zentraler Fortschritt der neueren Forschung besteht in der Einsicht, dass es nicht nur eine einzige Form des Judentums gab, von der andere Gruppen abgewichen wären. Vielmehr existierten unterschiedliche Interpretationen gemeinsamer religiöser Überlieferungen (Sanders 1992, 3–11; Grabbe 2010, 5–9). Die Tora bildete einen wichtigen Bezugspunkt, doch ihre Auslegung war umstritten: Welche Reinheitsgebote mussten eingehalten werden? Wer durfte den Tempelkult verwalten? Welche Stellung besaßen mündliche Traditionen? War die gegenwärtige politische Ordnung gottgewollt oder Ausdruck einer vorübergehenden Fremdherrschaft? Würde Gott die Toten auferwecken? War ein Messias zu erwarten, und welche Aufgaben sollte er übernehmen? Auf diese Fragen entwickelten verschiedene Gemeinschaften unterschiedliche Antworten.

Die häufig verwendete Einteilung in Pharisäer, Sadduzäer, Essener und Zeloten geht wesentlich auf Flavius Josephus zurück, der diese Gruppen für ein griechisch-römisches Publikum beschreibt und sich dabei der Sprache philosophischer Schulen bedient; seine Darstellung ist deshalb eine wichtige historische Quelle, aber keine vollständige Kartografie des antiken Judentums (Josephus, Antiquitates 18,11–25; Klawans 2012, 1–15; Taylor 2026, 209–224). Neben den von Josephus genannten Gruppen existierten priesterliche Traditionen, apokalyptische Kreise, jüdische Gemeinschaften in der Diaspora, Schriftgelehrte, lokale religiöse Vereinigungen und zahlreiche Menschen, die keiner bekannten Richtung eindeutig zugeordnet werden können. Das Judentum des Zweiten Tempels lässt sich daher besser als eine gemeinsame religiöse Lebenswelt mit unterschiedlichen Strömungen verstehen denn als ein geschlossenes konfessionelles System.

\section{Historischer Rahmen: Exil, Perser, Hellenismus und Rom}

Die Entstehung des Zweiten Tempel-Judentums ist ohne die Erfahrung des babylonischen Exils nicht zu verstehen. Mit der Eroberung Jerusalems durch die Babylonier im Jahr 587/586 v. Chr. wurden die politische Selbstständigkeit Judas und das religiöse Zentrum des Landes zerstört; Teile der Führungsschicht wurden deportiert, während andere Bewohner im Land verblieben (Grabbe 2010, 15–35). Die Zerstörung des Tempels bedeutete eine grundlegende Herausforderung für die religiöse Identität Israels: Wenn der Tempel als Ort der Gegenwart Gottes galt, stellte sich die Frage, wie diese Gegenwart ohne das Heiligtum gedacht werden konnte. Die exilische Literatur beantwortete diese Krise nicht einheitlich, entwickelte jedoch neue Deutungen von Bund, Gericht, Umkehr und göttlicher Treue.

Mit der Eroberung Babylons durch den persischen König Kyros II. im Jahr 539 v. Chr. änderten sich die politischen Voraussetzungen. Die persische Herrschaft ermöglichte jüdischen Exilierten die Rückkehr nach Juda und unterstützte die Wiederherstellung lokaler Heiligtümer. Die biblischen Bücher Esra und Nehemia stellen diese Entwicklung in den Zusammenhang einer religiösen Erneuerung (Kratz 2015, 10–18). Der Jerusalemer Tempel wurde um 515/516 v. Chr. fertiggestellt und bildete fortan den religiösen Mittelpunkt der persischen Provinz Yehud. Anders als zur Zeit der davidischen Monarchie existierte jedoch kein eigenständiges judäisches Königtum mehr: Die politische Herrschaft lag beim persischen Großkönig, während das Priestertum und lokale Verwaltungseliten innerhalb der jüdischen Gemeinschaft an Bedeutung gewannen. Die Reformen, die mit Esra und Nehemia verbunden werden, betonten die Verbindlichkeit der Tora, die religiöse Abgrenzung und die Ordnung des gemeinschaftlichen Lebens; ihre genaue historische Datierung und das Verhältnis ihrer Tätigkeiten werden in der Forschung diskutiert, und die biblischen Berichte sind deshalb nicht einfach als unmittelbare Chronik zu lesen, sondern zugleich als theologische Deutungen einer Neuordnung nach dem Exil (Kratz 2015, 18–25). Entscheidend ist die langfristige Entwicklung: Die religiöse Identität Israels beruhte zunehmend auf einer Verbindung von Tempel, Schrift, gemeinsamer Erinnerung und geregelter Lebenspraxis.

Mit den Eroberungen Alexanders des Großen ab 332 v. Chr. gelangte Juda unter den Einfluss der hellenistischen Welt. Nach Alexanders Tod wurde das Gebiet zunächst überwiegend von den ägyptischen Ptolemäern und später von den syrischen Seleukiden beherrscht (Grabbe 2010, 20–35). Die hellenistische Kultur brachte neue Formen der Bildung, Verwaltung, Sprache, Philosophie und städtischen Lebensweise mit sich; besonders das Griechische gewann als überregionale Verkehrssprache an Bedeutung. Jüdische Gemeinschaften begegneten diesen Entwicklungen in unterschiedlicher Weise: Hellenisierung bedeutete keineswegs zwangsläufig einen Verlust jüdischer Identität, und viele Juden verbanden die Verwendung der griechischen Sprache und die Teilnahme an hellenistischen Institutionen mit der Bindung an ihre religiösen Traditionen, während andere insbesondere in der Übernahme bestimmter kultureller Praktiken eine Bedrohung der überlieferten Ordnung sahen (Collins 2019, 440–455).

Unter Antiochos IV. Epiphanes verschärften sich die Konflikte; seine Eingriffe in die religiösen Verhältnisse Jerusalems und die Entweihung des Tempels führten ab 167 v. Chr. zum Makkabäeraufstand, und die Wiederweihe des Tempels 164 v. Chr. wurde später zum historischen Bezugspunkt des Chanukkafestes. Aus der makkabäischen Bewegung ging die hasmonäische Dynastie hervor, die zeitweise politische Selbstständigkeit erreichte. Die Hasmonäer verbanden priesterliche und politische Herrschaft, was jedoch neue Konflikte über die Legitimität des Hohepriestertums und die Ausübung königlicher Macht hervorrief; gerade diese Auseinandersetzungen trugen zur Ausbildung religiöser Gruppierungen bei, die unterschiedliche Vorstellungen über die richtige Ordnung des jüdischen Gemeinwesens entwickelten.

Im Jahr 63 v. Chr. eroberte der römische Feldherr Pompeius Jerusalem. Damit begann eine neue Phase politischer Abhängigkeit, die durch wechselnde Verwaltungsformen, lokale Herrschaftseliten und wiederkehrende Konflikte geprägt war (Sanders 1992, 15–25). Herodes der Große regierte von 37 bis 4 v. Chr. als von Rom gestützter König; seine Herrschaft war einerseits durch umfangreiche Bauprojekte und wirtschaftliche Entwicklung gekennzeichnet, andererseits durch politische Gewalt, dynastische Konflikte und die Abhängigkeit von Rom. Sein bedeutendstes religiöses Bauprojekt war der Ausbau des Jerusalemer Tempels, der um 20/19 v. Chr. begann; die monumentale Erweiterung der Tempelanlage machte Jerusalem zu einem eindrucksvollen religiösen Zentrum, das Pilger aus verschiedenen Regionen der jüdischen Diaspora anzog (Levine 1998, 53–68).

Nach dem Tod des Herodes wurde sein Herrschaftsgebiet unter mehreren Nachfolgern aufgeteilt; Judäa stand ab 6 n. Chr. zeitweise unter direkter römischer Verwaltung, und Pontius Pilatus amtierte von 26 bis 36 n. Chr. als Präfekt Judäas. Die römische Herrschaft verschärfte soziale und politische Spannungen: Steuerbelastungen, Konflikte um religiöse Symbole, Machtkämpfe zwischen lokalen Eliten und der römischen Verwaltung sowie unterschiedliche Erwartungen an die Wiederherstellung Israels trugen zu wiederholten Unruhen bei. Im Jahr 66 n. Chr. begann der große Jüdische Krieg gegen Rom; die römischen Truppen unter Titus eroberten Jerusalem im Jahr 70 n. Chr. und zerstörten den Tempel, und mit dem Fall Masadas 73/74 n. Chr. endete eine der letzten militärischen Widerstandsstellungen. Die Tempelzerstörung bedeutete eine tiefe Zäsur, aber nicht das Ende jüdischer Religionsgeschichte: Die Formen jüdischen Lebens, die bereits zuvor neben dem Tempelkult bestanden hatten, ermöglichten die weitere Entwicklung einer Religion ohne aktiven Opferkult.

\section{Sadduzäer und Pharisäer: Priesterliche Autorität und Tora-Auslegung}

Die Sadduzäer gehörten zu den bedeutenden religiösen und politischen Gruppen des späten Zweiten Tempel-Judentums. Sie standen besonders mit den priesterlichen und aristokratischen Kreisen Jerusalems in Verbindung, ohne dass sämtliche Priester als Sadduzäer gelten könnten (Klawans 2012, 15–30). Ihr religiöses Selbstverständnis war eng an die Autorität des geschriebenen Gesetzes und an die Ordnung des Tempelkults gebunden. Nach den neutestamentlichen Berichten und den Darstellungen des Josephus lehnten sie die Auferstehung der Toten ab; auch die pharisäische Vorstellung verbindlicher, über die schriftliche Tora hinausgehender Überlieferungen akzeptierten sie offenbar nicht in gleicher Weise (Josephus, Antiquitates 18,16–17; Klawans 2012, 30–45). Die häufig wiederholte Aussage, die Sadduzäer hätten grundsätzlich nicht an Engel geglaubt, bedarf einer gewissen Vorsicht: Sie beruht vor allem auf Apostelgeschichte 23,8 und lässt sich aus den übrigen Quellen nicht ohne Weiteres als umfassende Ablehnung aller biblischen Engelvorstellungen rekonstruieren; auch ihre genaue Haltung zur Autorität einzelner biblischer Bücher bleibt umstritten. Die Sadduzäer waren somit keine bloß politisch opportunistische Priestergruppe, sondern vertraten eigenständige theologische Positionen, die sich in Gesetzesverständnis, Tradition und Jenseitshoffnung deutlich von denen der Pharisäer unterschieden. Mit der Zerstörung des Jerusalemer Tempels verloren die sadduzäischen Eliten ihre wesentliche institutionelle Grundlage; als erkennbare religiöse Gruppierung verschwanden sie aus der späteren Überlieferung weitgehend.

Die Pharisäer gehören zu den bekanntesten Gruppen des antiken Judentums, nicht zuletzt aufgrund ihrer häufigen Erwähnung im Neuen Testament; die christliche Überlieferung hat ihr Bild allerdings über Jahrhunderte stark geprägt und teilweise verzerrt (Sanders 1992, 35–45). Historisch waren die Pharisäer eine religiöse Bewegung, die sich durch intensive Beschäftigung mit der Tora, durch Auslegungstraditionen und durch die Anwendung religiöser Vorschriften auf das alltägliche Leben auszeichnete. Neben dem geschriebenen Gesetz maßen sie überlieferten Auslegungen und Praktiken große Bedeutung bei; diese Traditionen ermöglichten es, die Gebote der Tora unter veränderten gesellschaftlichen Bedingungen anzuwenden. Ein wichtiges Merkmal war ihr Glaube an die Auferstehung der Toten und an eine zukünftige göttliche Gerechtigkeit; anders als die Sadduzäer vertraten sie damit eine Hoffnung, die über den Tod des Einzelnen hinausreichte (Josephus, Antiquitates 18,12–15; Sanders 1992, 45–55). Die verbreitete Bezeichnung als „volksnah“ ist jedoch nur eingeschränkt geeignet: Die Pharisäer verfügten nach Josephus über Einfluss in der Bevölkerung, waren aber zugleich in Auseinandersetzungen um religiöse Autorität und politische Macht eingebunden und können deshalb weder ausschließlich als Bewegung einfacher Gläubiger noch als homogene Gruppe von Schriftgelehrten beschrieben werden.

Für die Geschichte des Christentums sind die Pharisäer besonders wichtig, weil Jesus und seine ersten Anhänger in einer religiösen Welt lebten, in der pharisäische Gesetzesauslegung und religiöse Praxis weithin bekannt waren; viele neutestamentliche Streitgespräche sind vor diesem Hintergrund als innerjüdische Auseinandersetzungen zu verstehen. Nach der Tempelzerstörung gewannen Traditionen an Bedeutung, die mit dem pharisäischen Judentum verbunden waren; daraus entwickelte sich jedoch nicht unmittelbar und unverändert das rabbinische Judentum, sondern über Generationen erfolgte eine Neuorganisation religiöser Autorität, in der frühere Traditionen aufgenommen, umgeformt und neu interpretiert wurden. Die Frage nach Auferstehung und Jenseitshoffnung, die hier bereits im Unterschied von Sadduzäern und Pharisäern greifbar wird, wird in Kapitel 10 im weiteren frühjüdischen Horizont wieder aufgenommen.

\section{Essener, Qumran, Zeloten, Samaritaner und weitere Gemeinschaften}

Die Essener werden von Josephus, Philo und Plinius dem Älteren beschrieben; nach diesen Quellen zeichneten sie sich durch strenge religiöse Lebensführung, gemeinschaftliche Disziplin, Reinheitsvorschriften und eine besondere Distanz zu bestimmten gesellschaftlichen Praktiken aus (Grabbe 2010, 65–85). Einige Berichte betonen gemeinschaftlichen Besitz, geregelte Mahlzeiten und die Ablehnung individueller Bereicherung, doch ist zu beachten, dass die antiken Autoren unterschiedliche Formen essenischen Lebens beschreiben und nicht alle Essener notwendig in derselben Gemeinschaftsform lebten. Besondere Bedeutung erhielt die Essenerforschung durch die Entdeckung der Schriftrollen vom Toten Meer seit 1947; die Texte aus den Höhlen bei Qumran enthalten biblische Handschriften, Kommentare, Gemeinschaftsordnungen, liturgische Texte und apokalyptische Literatur (Grabbe 2010, 65–85). Die sogenannte Gemeinschaftsregel beschreibt eine religiöse Gruppe, die sich als besonders treue Gemeinschaft des Bundes verstand, zwischen den Angehörigen des Lichts und den Mächten der Finsternis unterschied und eine zukünftige göttliche Entscheidung erwartete.

Viele Forscher verbinden diese Gemeinschaft mit den Essenern oder mit einer ihnen verwandten religiösen Bewegung; eine vollständige Identifikation ist jedoch nicht gesichert, und die Qumran-Handschriften sind keine einheitliche Bibliothek einer einzigen theologischen Richtung, sondern enthalten Texte unterschiedlicher Entstehungszeiten und religiöser Herkunft. Besonders aufschlussreich ist die Kritik am Jerusalemer Tempelkult: Sie bedeutete nicht notwendig die grundsätzliche Ablehnung des Tempels als religiöser Institution, sondern richtete sich gegen ein als illegitim oder verunreinigt angesehenes Priestertum und gegen die bestehende religiöse Ordnung; die Gemeinschaft erwartete eine Erneuerung Israels und eine Wiederherstellung des rechten Gottesdienstes. Damit zeigt Qumran, dass selbst scharfe Kritik am Tempel aus einer tiefen Bindung an die religiöse Bedeutung des Heiligtums hervorgehen konnte. Die literarischen Eigenheiten der Qumran-Schriften werden in Kapitel 9 vertieft.

Neben den religiösen Auslegungsbewegungen entstanden Gruppen, die die römische Herrschaft militärisch bekämpften. Josephus beschreibt eine sogenannte vierte Philosophie, deren Entstehung er mit Judas dem Galiläer und den Auseinandersetzungen um den römischen Zensus von 6 n. Chr. verbindet; die Anhänger dieser Richtung betonten die alleinige Herrschaft Gottes und lehnten die politische Unterwerfung unter fremde Mächte ab (Josephus, Antiquitates 18,1–10; Grabbe 2010, 85–105). Die später als Zeloten bezeichneten Akteure spielten im Jüdischen Krieg eine wichtige Rolle; die Sikarier, deren Bezeichnung auf kurze Dolche zurückgeht, waren besonders durch politische Gewalt und gezielte Anschläge bekannt. Eine Gleichsetzung aller antirömischen Gruppen ist jedoch problematisch, weil die Begriffe nicht durchgängig dieselbe Bewegung bezeichnen und ihre Entstehungsgeschichte sich nicht vollständig rekonstruieren lässt; auch die religiösen Motive waren vielfältig, da politische Unabhängigkeit, soziale Konflikte, priesterliche Rivalitäten und die Hoffnung auf göttliches Eingreifen ineinandergreifen konnten. Die militärische Niederlage des jüdischen Aufstands führte zum Zusammenbruch zahlreicher politischer Organisationsformen, doch die religiösen Hoffnungen auf die Befreiung Israels verschwanden dadurch nicht.

Die religiöse Landschaft der Epoche beschränkte sich nicht auf die Gruppen, die Josephus beschreibt. Die Samaritaner bildeten eine eigenständige, mit dem Israelitentum historisch verwandte Religionsgemeinschaft; sie verehrten den Gott Israels, erkannten den Pentateuch als maßgebliche Schrift an und betrachteten den Berg Garizim als den rechtmäßigen Ort des Heiligtums. Ihr Verhältnis zum Jerusalemer Judentum war von gemeinsamen Überlieferungen, aber auch von Auseinandersetzungen über Kultort, Abstammung und religiöse Legitimität geprägt; diese Spannung verschärfte sich in hasmonäischer Zeit militärisch, als der Hasmonäer Johannes Hyrkanos I. um 110 v. Chr. das samaritanische Heiligtum auf dem Garizim sowie die benachbarte Tempel- und Priesterstadt zerstörte (Josephus, Antiquitates 13,9,1; archäologisch bestätigt durch die Ausgrabungen Yitzhak Magens, die eine Zerstörung um 111/110 v. Chr. nahelegen). Die Zerstörung war kein einmaliger „Sonderfall“, sondern markiert den gewaltsamen Höhepunkt eines langen, wechselseitigen Abgrenzungsprozesses; sie zeigt, wie eng religiöse Differenz, priesterliche Konkurrenz und politisch-militärische Machtinteressen ineinandergriffen (vgl. Hensel 2014, 475–479). Der Begriff „judäisch-samaritanischer Konflikt“ ist dort präziser, wo die historischen Akteure aus Judäa und Samaria einander gegenüberstehen.

In Ägypten erwähnt Philo außerdem die Therapeuten, eine Gemeinschaft, deren Lebensweise durch religiöse Kontemplation, Schriftauslegung und asketische Praktiken gekennzeichnet gewesen sein soll; ihre genaue historische Einordnung bleibt umstritten (Philo, De vita contemplativa; Grabbe 2010, 115–125). Auch außerhalb solcher Gruppen bestanden vielfältige Formen jüdischen Lebens: Priester, Schriftgelehrte, Händler, Handwerker, Bauern und jüdische Gemeinschaften in den Städten des Mittelmeerraums unterschieden sich in ihren Lebensbedingungen und religiösen Ausdrucksformen. Die Vielfalt des Zweiten Tempel-Judentums ist deshalb nicht allein eine Vielfalt theologischer Schulen, sondern ebenso eine Vielfalt sozialer Milieus, regionaler Traditionen und kultureller Lebensweisen.

\section{Tempel, Synagoge und Tora}

Trotz aller religiösen Vielfalt blieb der Jerusalemer Tempel bis zu seiner Zerstörung im Jahr 70 n. Chr. der wichtigste gemeinsame Bezugspunkt des jüdischen Lebens (Sanders 1992, 75–95; Levine 1998, 53–68). Der Tempel war nicht lediglich ein Gebäude, in dem Opfer dargebracht wurden, sondern verkörperte die besondere Beziehung zwischen Gott und Israel, die Erinnerung an den Bund und die Hoffnung auf die fortdauernde Gegenwart Gottes inmitten seines Volkes. Im Mittelpunkt stand der Opferkult, dessen Durchführung dem Priestertum vorbehalten war; zu den regelmäßigen Kulthandlungen gehörten die täglichen Opfer sowie besondere Opfer an Sabbaten und Festtagen, daneben wurden persönliche Opfer aus unterschiedlichen Anlässen dargebracht, etwa im Zusammenhang mit Dank, Reinigung oder Sühne. Die religiöse Bedeutung des Tempels zeigte sich besonders an den drei Wallfahrtsfesten Pessach, Schawuot und Sukkot: Pessach erinnerte an die Befreiung Israels aus Ägypten und verband historische Erinnerung mit dem Bekenntnis zu Gottes rettendem Handeln; Schawuot war ursprünglich ein Erntefest und gewann später zunehmend Bedeutung im Zusammenhang mit der Gabe der Tora; Sukkot erinnerte an die Wüstenwanderung und verband diese Erinnerung mit Dankbarkeit für die Ernte.

Der Tempel erfüllte zugleich wirtschaftliche und gesellschaftliche Funktionen: Die Versorgung des Opferkults, die Tempelabgaben, der Pilgerverkehr und die Verwaltung des Heiligtums schufen ein komplexes institutionelles Gefüge, und die priesterlichen Familien verfügten dadurch über religiöse Autorität und teilweise beträchtlichen politischen Einfluss. Besonders das Hohepriesteramt war während der hellenistischen und römischen Herrschaft Gegenstand politischer Auseinandersetzungen. Dennoch wäre es falsch, den Tempel ausschließlich als Herrschaftsinstrument zu verstehen: Für die Mehrheit der jüdischen Bevölkerung stellte er einen Ort gemeinsamer religiöser Identität dar, dessen Bedeutung die jeweiligen Konflikte um die Verwaltung des Kultes überstieg.

Parallel zum Jerusalemer Tempel entwickelte sich die Synagoge als eigenständiger Ort jüdischer Gemeinschaft (Levine 2000, 1–45). Ihre genaue Entstehungsgeschichte ist historisch nicht vollständig geklärt; häufig wird ihre Entwicklung mit der Erfahrung des Exils und den Bedürfnissen jüdischer Gemeinschaften außerhalb Jerusalems verbunden, doch eine unmittelbare Gründung der Synagoge während des babylonischen Exils ist nicht nachweisbar. Archäologische und schriftliche Zeugnisse belegen synagogenartige Versammlungsstätten in der hellenistischen und römischen Zeit; besonders in Ägypten begegnet die griechische Bezeichnung proseuchē, die einen Gebets- oder Versammlungsort bezeichnet. Die Synagoge erfüllte mehrere Funktionen: Sie diente der gemeinschaftlichen Zusammenkunft, dem Lesen und Auslegen der Tora, der religiösen Unterweisung und zunehmend auch dem Gebet. Damit unterscheidet sie sich grundlegend vom Jerusalemer Tempel: Während der Tempel der zentrale Ort des Opferkults war, konnte eine Synagoge grundsätzlich überall dort entstehen, wo sich eine jüdische Gemeinde organisierte.

Die beiden Institutionen standen zunächst nicht in Konkurrenz, sondern ergänzten sich; Juden, die regelmäßig eine Synagoge besuchten, konnten zugleich die Bedeutung des Jerusalemer Tempels anerkennen, die Tempelsteuer entrichten und an Wallfahrtsfesten teilnehmen (Cohen 1999, 298–325). Erst nach der Zerstörung des Tempels gewann die Synagoge eine neue religiöse Bedeutung und wurde schrittweise zu einem zentralen Ort jüdischer Kontinuität; dieser Wandel vollzog sich jedoch regional unterschiedlich und nicht unmittelbar im gesamten Judentum. Die Entwicklung von Tora, Synagoge und Gebet ist deshalb nicht als geradlinige Vorbereitung auf eine Religion ohne Tempel zu verstehen: Diese Formen existierten bereits parallel zum Tempelkult und hatten eigene Funktionen; nach 70 n. Chr. veränderten sich ihre Gewichtung und ihre institutionellen Rahmenbedingungen, ohne dass daraus ein einfacher Ersatz des Tempels wurde.

Neben dem Tempel bildete die Tora das zentrale Fundament jüdischer Identität (Sanders 1992, 95–105). Der Begriff „Tora“ bedeutet zunächst Weisung oder Unterweisung und bezeichnet im engeren Sinne die fünf Bücher Mose, in einem weiteren Verständnis die göttliche Orientierung des gesamten religiösen Lebens. Die Tora regelte nicht lediglich kultische Handlungen, sondern betraf auch soziale Beziehungen, Eigentum, Rechtsprechung, Familienleben, Speisevorschriften und den Umgang mit gesellschaftlich benachteiligten Menschen; gerade diese umfassende Bedeutung ermöglichte eine religiöse Lebensführung auch fern vom Jerusalemer Tempel. Die Auslegung der Tora war allerdings nicht einheitlich: Unterschiedliche Gemeinschaften entwickelten unterschiedliche Vorstellungen darüber, wie die Gebote unter den Bedingungen der Gegenwart anzuwenden seien (Grabbe 2010, 125–135). Fragen der rituellen Reinheit, des Sabbats, der Ehe oder des Kalenders konnten zu ernsthaften religiösen Konflikten führen; die Qumran-Gemeinschaft vertrat beispielsweise eigene Kalendertraditionen, die sie vom Jerusalemer Tempelkult unterschieden. Die religiöse Beschäftigung mit der Tora war damit nicht nur eine Frage der Bewahrung älterer Vorschriften, sondern zugleich ein kreativer Prozess der Interpretation.

\section{Religiöse Lebenspraxis: Gebet, Reinheit und Alltag}

Das religiöse Leben des Zweiten Tempel-Judentums erschöpfte sich keineswegs in Tempelkult und Schriftauslegung. Gebet, Fasten, Almosen, Speisegebote, Sabbatruhe und Reinheitspraktiken prägten den Alltag in unterschiedlichem Umfang; wie intensiv einzelne Vorschriften befolgt wurden, hing unter anderem von regionalen, sozialen und religiösen Bedingungen ab (Sanders 1992, 105–125). Rituelle Reinheit besaß eine besondere Bedeutung, weil sie bestimmte Formen der Teilnahme am Kult ermöglichte; archäologische Funde zahlreicher ritueller Tauchbäder (Mikwaot) belegen die Verbreitung entsprechender Praktiken, insbesondere in Judäa während der späten Zweiten Tempelzeit. Dabei darf rituelle Unreinheit nicht mit moralischer Schuld gleichgesetzt werden: Sie konnte beispielsweise durch normale körperliche Vorgänge oder den Kontakt mit Toten entstehen und war Teil einer religiösen Ordnung, die zwischen unterschiedlichen Zuständen und Handlungen unterschied; zugleich entwickelten sich Vorstellungen, die kultische Reinheit mit ethischer Verantwortung verbanden. Die Propheten hatten bereits betont, dass Opfer ohne Gerechtigkeit und Treue zum Bund ihren Sinn verfehlen; diese Tradition wirkte im Frühjudentum weiter, sodass Frömmigkeit nicht allein durch kultische Handlungen bestimmt wurde, sondern ebenso durch verantwortliches Handeln gegenüber anderen Menschen.

\section{Schriftwerdung, Kanon und literarische Überlieferung}

Die Zeit des Zweiten Tempels war eine außerordentlich produktive Phase jüdischer Literaturgeschichte (Grabbe 2010, 135–155; Collins 2019, 437–459). Viele Schriften der späteren Hebräischen Bibel erhielten in dieser Epoche ihre endgültige oder annähernd endgültige Gestalt; zugleich entstanden zahlreiche weitere religiöse Texte, die biblische Überlieferungen fortführten, auslegten oder neu gestalteten. Die Herausbildung einer Sammlung heiliger Schriften war jedoch kein einmaliger Vorgang: Die Tora besaß bereits relativ früh besondere Autorität, während die Grenzen anderer Schriftgruppen länger in Bewegung blieben. Die Qumran-Handschriften zeigen, dass verschiedene Textfassungen nebeneinander existierten; neben Vorformen des späteren masoretischen Textes begegnen Handschriften, deren Wortlaut anderen Texttraditionen nähersteht. Auch die griechische Übersetzung jüdischer Schriften, die später unter dem Namen Septuaginta bekannt wurde, bezeugt eine vielfältige Überlieferungsgeschichte. Von einer vollständig abgeschlossenen Kanonisierung der gesamten Hebräischen Bibel während der Zweiten Tempelzeit kann daher nicht ohne Einschränkungen gesprochen werden. Die spätere Unterscheidung zwischen kanonischen, apokryphen und pseudepigraphen Schriften darf nicht unkritisch auf die damalige religiöse Wirklichkeit übertragen werden; Texte, die später außerhalb eines bestimmten Kanons standen, konnten in einzelnen frühjüdischen Gemeinschaften erhebliche Autorität besitzen.

Zu den wichtigen literarischen Zeugnissen der Epoche gehören die Bücher Tobit, Judit, Jesus Sirach, die Weisheit Salomos und die Makkabäerbücher. Sie behandeln unterschiedliche Fragen der religiösen Existenz: die Bewahrung des Glaubens in der Fremde, die Bedeutung der Weisheit, den Umgang mit politischer Unterdrückung, das Martyrium und die Hoffnung auf göttliche Gerechtigkeit. Das Buch Tobit verbindet die Erzählung einer frommen jüdischen Familie in der Diaspora mit Weisheitslehren, Gebet und dem Wirken eines Engels; das Buch Judit erzählt von einer Frau, die durch Mut und Klugheit zur Rettung ihres Volkes beiträgt, und verbindet religiöse Treue mit einer dramatischen Erzählung politischer Bedrohung; Jesus Sirach entwirft eine umfassende Weisheitslehre, in der die Weisheitstraditionen Israels ausdrücklich mit der Tora verbunden werden; die Weisheit Salomos zeigt besonders deutlich die Begegnung jüdischer Theologie mit hellenistischer Begrifflichkeit und behandelt Fragen der Gerechtigkeit, Unsterblichkeit, Weisheit und göttlichen Weltordnung. Die Makkabäerbücher sind für die Geschichte der religiösen Verfolgung und des jüdischen Widerstands von besonderem Wert; vor allem das zweite Makkabäerbuch enthält ausgeprägte Vorstellungen von Auferstehung und göttlichem Gericht. Diese Schriften wurden in verschiedenen jüdischen und christlichen Traditionen unterschiedlich bewertet: Während einige in katholischen und orthodoxen Bibelkanones aufgenommen wurden, gehören sie nicht zum jüdischen Tanach; der Sammelbegriff „Apokryphen“ ist deshalb stets im Zusammenhang der jeweiligen Kanontradition zu verstehen.

Die sogenannte pseudepigraphe Literatur bildet einen besonders aufschlussreichen Teil der frühjüdischen Überlieferung. Als Pseudepigraphen werden Schriften bezeichnet, die unter dem Namen einer bekannten Gestalt der religiösen Vergangenheit auftreten, obwohl sie aus späterer Zeit stammen; diese Praxis muss im Zusammenhang antiker literarischer Autoritätskonzepte verstanden werden und darf nicht ohne Weiteres mit moderner literarischer Täuschung gleichgesetzt werden. Zu den bedeutendsten Werken gehören das äthiopische Henochbuch, das Jubiläenbuch, die Psalmen Salomos und verschiedene Testamentenliteraturen. Das Henochbuch besteht aus mehreren literarischen Einheiten unterschiedlicher Entstehungszeit; es verbindet Visionen des himmlischen Bereichs mit Vorstellungen über Engel, Gericht, kosmische Ordnung und die Zukunft der Gerechten, wobei die Erzählung von Engeln, die ihre himmlische Ordnung verlassen und dadurch Unheil in die Welt bringen, eine zentrale Rolle spielt und ältere biblische Motive weiterentwickelt. Das Jubiläenbuch erzählt die Geschichte von der Schöpfung bis zur Zeit des Mose neu, gliedert sie in festgelegte Zeitabschnitte und betont die Bedeutung des göttlichen Gesetzes, des Kalenders und der religiösen Reinheit. Die Psalmen Salomos verbinden die Kritik an politischen Herrschern mit der Hoffnung auf einen gerechten davidischen König und gehören damit zu den wichtigen Zeugnissen frühjüdischer Messiaserwartung. Die Pseudepigraphen zeigen, wie ältere biblische Überlieferungen unter veränderten historischen Bedingungen fortgeschrieben wurden; ihre Autoren verstanden die religiöse Vergangenheit nicht als abgeschlossene Geschichte, sondern als Quelle neuer Einsichten über Gegenwart und Zukunft. Allerdings sind diese Schriften weder durchweg apokalyptisch noch sämtlich eindeutig vorchristlich; einzelne Werke oder Textschichten besitzen eine komplexe Überlieferungsgeschichte und können spätere jüdische oder christliche Bearbeitungen enthalten.

\section{Qumran, Weisheit, Logos und hellenistische Philosophie}

Die Entdeckung der Handschriften vom Toten Meer gehört zu den wichtigsten Ereignissen der modernen Bibelwissenschaft. Seit 1947 wurden in den Höhlen bei Qumran zahlreiche Handschriftenfragmente gefunden, die einen außergewöhnlichen Einblick in die jüdische Literatur der letzten Jahrhunderte vor der Zerstörung des Zweiten Tempels eröffnen. Die Sammlung umfasst biblische Texte, Auslegungen, Gemeinschaftsordnungen, liturgische Werke, Kalendertexte, Weisheitsliteratur und apokalyptische Schriften. Von besonderer Bedeutung ist die Gemeinschaftsregel (1QS), die die Aufnahme neuer Mitglieder, die Ordnung des gemeinschaftlichen Lebens und die religiösen Pflichten der Angehörigen beschreibt; die Damaskusschrift enthält Regelungen und Ermahnungen, die eine Gemeinschaft als Trägerin des wahren Bundes verstehen lassen; die Kriegsrolle (1QM) entwirft die Vorstellung eines endzeitlichen Kampfes zwischen den Angehörigen des Lichts und den Mächten der Finsternis; die sogenannten Peschärim bieten Auslegungen prophetischer Texte, in denen die alten Verheißungen auf die eigene Gemeinschaft und ihre historischen Erfahrungen bezogen werden. Damit bezeugen die Qumran-Texte eine charakteristische Form religiöser Selbstdeutung: Eine Gemeinschaft versteht sich als besonderer Träger göttlicher Erkenntnis und als Teil eines heilsgeschichtlichen Geschehens, dessen Vollendung unmittelbar bevorsteht. Für die neutestamentliche Forschung sind diese Texte besonders wichtig, weil sie zeigen, dass Begriffe und Motive wie Licht und Finsternis, Bund, Geist, endzeitlicher Kampf und messianische Erwartung bereits innerhalb des Judentums vielfältig entwickelt waren; daraus folgt allerdings keine unmittelbare Abhängigkeit Jesu oder der frühen Christen von der Qumran-Gemeinschaft, vielmehr belegen die Texte einen gemeinsamen religiösen und sprachlichen Horizont. Die bereits in Kapitel 5 dargestellte Verbindung zwischen Essenern und Qumran wird hier also nicht erneut entfaltet, sondern durch die literarischen Zeugnisse ergänzt.

In der neueren Forschung wird diskutiert, ob sich hinter der Henochliteratur eine eigenständige „enochische“ Tradition gegenüber einer „zadokitisch“-priesterlichen Überlieferung profilieren lässt. Gabriele Boccaccini hat diese Unterscheidung stark gemacht und die These vertreten, die Wurzeln der Qumran-Gemeinde lägen weniger in einer zadokitischen „Reaktion“ als vielmehr in der enochischen Tradition (Boccaccini 1998; Boccaccini 2002). Das Modell ist heuristisch fruchtbar, bleibt jedoch umstritten, weil die Qumran-Texte keine saubere Trennung zweier „Konfessionen“ zeigen, sondern vielfältige Überschneidungen; die Gemeinschaft ist zugleich zadokitisch-priesterlich und apokalyptisch geprägt, und 1 Henoch ist eine Sammlung unterschiedlicher Textschichten, die nicht ohne Weiteres auf ein einziges „Judentum“ zurückgeführt werden kann (vgl. die kritische Diskussion bei Albani 2007, 85–101).

Ein weiterer Schwerpunkt der Epoche liegt in der Auseinandersetzung jüdischer Traditionen mit der griechischen Philosophie, die besonders im Werk des jüdischen Philosophen Philo von Alexandria deutlich wird, der ungefähr von 20 v. Chr. bis 50 n. Chr. lebte. Philo war tief in der jüdischen Schrifttradition verwurzelt und zugleich mit philosophischen Strömungen seiner Zeit vertraut, insbesondere mit platonischen und stoischen Vorstellungen. Sein Anliegen bestand darin, den philosophischen Gehalt der Tora sichtbar zu machen; dazu verwendete er vielfach die allegorische Schriftauslegung, bei der biblische Erzählungen neben ihrer wörtlichen Bedeutung auch als Ausdruck geistiger und ethischer Wahrheiten verstanden werden. Eine zentrale Stellung nimmt in seinem Denken der Logos ein: Der Begriff bezeichnet in unterschiedlichen Zusammenhängen die göttliche Vernunft, das Wort oder die vermittelnde Ordnung, durch die Gott mit der geschaffenen Welt in Beziehung tritt; der Logos ist bei Philo allerdings nicht einfach eine von Gott getrennte zweite Gottheit, und die entsprechenden Aussagen sind vielschichtig und müssen im Zusammenhang seines monotheistischen Denkens verstanden werden. Für die Geschichte christlicher Theologie ist die Logos-Tradition bedeutsam, weil das Johannesevangelium seine Christologie ebenfalls mit diesem Begriff entfaltet; eine direkte Übernahme aus Philo lässt sich daraus jedoch nicht ableiten, vielmehr stehen beide in einem weiteren jüdischen und hellenistischen Diskussionszusammenhang. Neben Philo zeigt auch die Weisheit Salomos, dass jüdische Theologie hellenistische Begriffe aufnehmen konnte, ohne ihre Bindung an den Gott Israels aufzugeben.

\section{Apokalyptik, Auferstehung und messianische Erwartungen}

Die Apokalyptik gehört zu den prägenden religiösen Ausdrucksformen des Frühjudentums. Der Begriff leitet sich vom griechischen apokalypsis, Offenbarung, ab; gemeint sind literarische und theologische Traditionen, in denen verborgene Zusammenhänge der Welt und der Geschichte durch göttliche Offenbarung sichtbar werden. Apokalyptische Texte arbeiten häufig mit Visionen, Engeldeutungen, symbolischen Bildern und der Gegenüberstellung gegenwärtiger Bedrängnis und zukünftiger Erlösung. Besonders einflussreich ist das Danielbuch, dessen Visionen die Herrschaft irdischer Reiche als vorübergehend darstellen und die endgültige Durchsetzung göttlicher Herrschaft erwarten lassen. Apokalyptik entsteht häufig in Situationen politischer Bedrohung und religiöser Verunsicherung; sie bietet eine Deutung der Geschichte, in der die scheinbare Macht der Unterdrücker nicht das letzte Wort behält. Dieser Zusammenhang mit politischen Krisen ist allerdings nicht als allgemeine Entstehungserklärung sämtlicher apokalyptischer Literatur zu verstehen: Apokalyptische Texte können auch ohne akute Bedrohung entstehen und umfassen ebenso Visionen des Himmels, der kosmischen Ordnung und verborgener göttlicher Wirklichkeit. Der gegenwärtigen Weltordnung wird eine zukünftige göttliche Ordnung gegenübergestellt, wobei die Hoffnung auf eine grundlegende Transformation der Geschichte mit Vorstellungen von Gericht, Auferstehung und einer neuen Schöpfung verbunden sein kann. Für das frühe Christentum war dieses Denken von grundlegender Bedeutung: Vorstellungen vom kommenden Gottesreich, vom Menschensohn, von der Auferstehung und vom endzeitlichen Gericht stehen in einem frühjüdisch-apokalyptischen Zusammenhang.

Die Frage nach dem Schicksal des Menschen nach dem Tod erhielt im Frühjudentum neue Bedeutung (Sanders 1992, 125–145). Ältere biblische Vorstellungen beschreiben den Tod häufig als Eintritt in die Scheol, den Bereich der Toten, in dem die Möglichkeiten menschlichen Lebens begrenzt sind; im Laufe der Zeit entwickelten sich daneben deutlichere Hoffnungen auf eine zukünftige Überwindung des Todes. Daniel 12,2 spricht davon, dass viele der im Staub Schlafenden erwachen werden, die einen zum ewigen Leben, die anderen zu Schmach und Abscheu; auch das zweite Makkabäerbuch enthält ausgeprägte Auferstehungshoffnungen, insbesondere im Zusammenhang mit dem Martyrium. Diese Vorstellungen beantworten eine grundlegende religiöse Frage: Wie kann Gottes Gerechtigkeit bestehen, wenn gerechte Menschen leiden und sterben, während ihre Unterdrücker scheinbar ungestraft bleiben? Die Auferstehungshoffnung verlagert die Erwartung göttlicher Gerechtigkeit über die Grenzen des gegenwärtigen Lebens hinaus. Allerdings existierte keine einheitliche frühjüdische Lehre über das Leben nach dem Tod; neben Vorstellungen körperlicher Auferstehung begegnen unterschiedliche Aussagen über die Unsterblichkeit der Seele, die Existenz der Gerechten bei Gott und eine zukünftige Erneuerung der Welt. Gerade die bereits in Kapitel 4 dargestellten Unterschiede zwischen Sadduzäern und Pharisäern zeigen, dass solche Fragen innerhalb des Judentums intensiv diskutiert wurden.

Die Vorstellung eines Messias war im Judentum des Zweiten Tempels keineswegs einheitlich. Der Begriff „Messias“ geht auf das hebräische maschiach, Gesalbter, zurück und bezeichnet ursprünglich Personen, die durch Salbung zu einem besonderen Amt bestimmt wurden, insbesondere Könige und Priester. In der frühjüdischen Literatur entwickelten sich unterschiedliche Erwartungen an zukünftige Heilsgestalten: Eine wichtige Tradition erwartete einen davidischen König, der Israel erneuern und eine gerechte Herrschaft errichten würde; andere Texte verbinden die Hoffnung auf Heil mit einer priesterlichen Gestalt oder mit mehreren endzeitlichen Beauftragten Gottes. In einzelnen Qumran-Texten begegnen Vorstellungen von messianischen Gestalten mit königlichen und priesterlichen Funktionen; daneben finden sich apokalyptische Visionen himmlischer Gestalten, beispielsweise in den Bilderreden des Henochbuches, wobei die Datierung dieser Textschicht umstritten ist und sie für die Rekonstruktion vorchristlicher Messiasvorstellungen nur mit Vorsicht herangezogen werden kann. Nicht sämtliche frühjüdischen Gemeinschaften erwarteten einen Messias im gleichen Sinne; auch bestand keine allgemein verbindliche Vorstellung darüber, ob der Messias ein politischer Befreier, ein Priester, ein himmlischer Richter oder eine andere Heilsgestalt sein würde. Die spätere christliche Identifikation Jesu als Messias entstand innerhalb dieser vielfältigen religiösen Erwartungswelt.

\section{Diaspora, Septuaginta und hellenistisch-römische Welt}

Ein wesentlicher Teil der jüdischen Bevölkerung lebte bereits während der Zweiten Tempelzeit außerhalb Judäas (Grabbe 2010, 185–205). Jüdische Gemeinschaften bestanden in Babylonien, Ägypten, Syrien, Kleinasien, Griechenland und Rom; ihre Entstehung beruhte auf unterschiedlichen historischen Entwicklungen, darunter Deportationen, freiwillige Migration, Handel und militärische Ansiedlungen. Besonders Alexandria entwickelte sich zu einem bedeutenden Zentrum jüdischen Lebens: Die dortigen Juden nahmen am kulturellen Leben der hellenistischen Stadt teil und bewahrten zugleich ihre religiösen Traditionen. Die Diaspora war damit nicht bloß eine vorübergehende Folge politischer Vertreibung, sondern wurde zu einer dauerhaften Form jüdischer Existenz. Für die Religionsgeschichte ist dies von grundlegender Bedeutung: Jüdische Identität konnte nun über die Zugehörigkeit zu einer lokalen Gemeinschaft, die religiöse Praxis und die gemeinsame Schrifttradition bewahrt werden, ohne dass ein dauerhafter Aufenthalt im Land Israel erforderlich war. Trotzdem blieb Jerusalem für zahlreiche Diasporagemeinden ein religiöses Zentrum; Pilgerreisen, Geldsendungen für den Tempel und gemeinsame Festtraditionen verbanden die verstreuten Gemeinschaften.

Die Septuaginta gehört zu den bedeutendsten Ergebnissen jüdischer Kulturgeschichte in der hellenistischen Welt. Unter diesem Namen wird die griechische Übersetzung jüdischer heiliger Schriften verstanden, deren Anfänge in Alexandria gewöhnlich in das 3. Jahrhundert v. Chr. datiert werden. Der Überlieferung nach wurde die Tora von zweiundsiebzig jüdischen Gelehrten übersetzt; diese Erzählung findet sich im Aristeasbrief und dient zugleich der religiösen Legitimation der griechischen Übersetzung. Historisch entwickelte sich die Septuaginta über einen längeren Zeitraum: Die einzelnen Bücher wurden von unterschiedlichen Übersetzern bearbeitet und spiegeln unterschiedliche sprachliche und theologische Voraussetzungen wider. Die griechische Übersetzung machte die jüdischen Schriften für Menschen zugänglich, die Hebräisch nicht oder nur eingeschränkt beherrschten, und eröffnete zugleich neue Möglichkeiten der Interpretation, weil Begriffe der hebräischen Bibel in eine andere Sprache und damit in einen teilweise anderen begrifflichen Horizont übertragen werden mussten. Die Septuaginta war deshalb nicht nur eine sprachliche Vermittlung, sondern auch ein bedeutender Ort theologischer Auslegung; für das frühe Christentum wurde sie besonders wichtig, weil viele neutestamentliche Schriftzitate auf griechischen Textfassungen beruhen. Die kanongeschichtlichen Voraussetzungen dieses Überlieferungsprozesses wurden bereits in Kapitel 8 skizziert und müssen hier nicht erneut entfaltet werden.

Die Begegnung mit der hellenistischen Welt führte zu unterschiedlichen Formen jüdischer Selbstbehauptung (Grabbe 2010, 195–205). Einige jüdische Autoren bemühten sich darum, die Übereinstimmung jüdischer Weisheit mit philosophischen Traditionen der griechischen Welt hervorzuheben; andere betonten die Unterscheidung zwischen der Verehrung des Gottes Israels und den religiösen Praktiken ihrer Umwelt. Sabbat, Beschneidung, Speisegebote und die Ablehnung fremder Götter konnten zu sichtbaren Merkmalen jüdischer Zugehörigkeit werden; diese Praktiken besaßen jedoch nicht ausschließlich die Funktion einer Abgrenzung, sondern strukturierten ebenso das gemeinschaftliche Leben und vermittelten religiöse Erinnerung. Die Begriffe „Assimilation“ und „Abgrenzung“ beschreiben daher nur zwei Pole eines komplexeren Verhältnisses: Jüdische Gemeinschaften konnten griechische Sprache, Bildung und Verwaltungsformen übernehmen und zugleich an ihrem besonderen religiösen Selbstverständnis festhalten. Die kulturelle Begegnung führte nicht notwendig zu einer Auflösung der jüdischen Tradition, sondern konnte ebenso neue Möglichkeiten ihrer theologischen Entfaltung schaffen.

\section{Frühes Christentum, rabbinisches Judentum und Schlussbetrachtung}

Jesus von Nazaret lebte und wirkte innerhalb der religiösen und gesellschaftlichen Welt des Judentums des Zweiten Tempels (Sanders 1992, 145–165). Seine Verkündigung vom Reich Gottes, seine Schriftauslegung, seine Auseinandersetzungen um den Sabbat, seine Stellung zum Tempel und seine Hoffnung auf Gottes zukünftiges Handeln sind ohne diesen Hintergrund nicht verständlich. Jesus war kein Begründer einer bereits von Anfang an eigenständigen, vom Judentum getrennten Religionsgemeinschaft; seine öffentliche Wirksamkeit richtete sich zunächst innerhalb Israels an jüdische Menschen und griff auf gemeinsame biblische Überlieferungen zurück. Auch die Auseinandersetzungen mit Pharisäern und anderen religiösen Autoritäten sind vor diesem Hintergrund zu betrachten und lassen sich nicht angemessen als Gegensatz zwischen einer angeblich starren jüdischen Gesetzesreligion und einer völlig neuen christlichen Religion darstellen. Vielmehr gehören sie zu den innerjüdischen Debatten über die Auslegung der Tora, die Bedeutung religiöser Reinheit, die Stellung des Menschen vor Gott und die Verwirklichung göttlicher Gerechtigkeit; gerade das Neue Testament selbst bezeugt diese Vielfalt jüdischer Positionen.

Auch Paulus von Tarsus lässt sich nur innerhalb der jüdischen Welt seiner Zeit angemessen verstehen. Nach seiner eigenen Darstellung gehörte er zur jüdischen Tradition und verstand sich in enger Verbindung mit dem Gesetz und den Überlieferungen seiner Vorfahren (Phil 3,5–6; Gal 1,13–14). Seine Verkündigung von Jesus Christus als dem auferstandenen Messias führte jedoch zu neuen Fragen über die Aufnahme nichtjüdischer Menschen in die Gemeinschaft der Christusgläubigen: Müssen nichtjüdische Anhänger Jesu beschnitten werden? Welche Bedeutung besitzen die Speisegebote? Wie verhalten sich die Verheißungen an Israel zur Einbeziehung der Völker? Diese Auseinandersetzungen zeigen, dass die Entstehung des frühen Christentums nicht einfach die Ablösung des Judentums bedeutete, sondern sich zunächst als Streit über die Reichweite und die Auslegung jüdischer Verheißungen vollzog. Paulus argumentierte dabei mit der jüdischen Schrifttradition, bediente sich jedoch zugleich der griechischen Sprache und rhetorischer Formen seiner Umwelt; sein Denken verbindet somit unterschiedliche Elemente der jüdischen Diasporakultur.

Die Zerstörung des Jerusalemer Tempels im Jahr 70 n. Chr. stellte eine der tiefgreifendsten Krisen der jüdischen Geschichte dar (Grabbe 2010, 205–225). Mit dem Heiligtum verlor das Judentum den zentralen Ort seines Opferkults; zugleich wurden die politischen und wirtschaftlichen Grundlagen der Jerusalemer Priesteraristokratie zerstört. Die religiöse Reaktion auf diese Katastrophe verlief keineswegs einheitlich: Apokalyptische Schriften verarbeiteten den Verlust Jerusalems als Ausdruck einer tiefen geschichtlichen Krise und fragten nach Gottes Gerechtigkeit, während andere jüdische Traditionen die fortdauernde Bedeutung der Tora, des Gebets, der religiösen Praxis und der gemeinschaftlichen Auslegung betonten. In diesem Zusammenhang gewann die rabbinische Bewegung langfristig an Bedeutung. Die Tradition verbindet einen frühen Mittelpunkt rabbinischer Gelehrsamkeit mit Jabne beziehungsweise Jamnia; die ältere Vorstellung, dort habe nach 70 n. Chr. ein verbindliches Konzil den jüdischen Bibelkanon festgelegt, gilt jedoch als historisch nicht haltbar. Ebenso wenig entstand das rabbinische Judentum unmittelbar durch einen einmaligen organisatorischen Beschluss; vielmehr entwickelte sich über mehrere Generationen eine neue Form religiöser Autorität, in deren Zentrum Gelehrte und ihre Überlieferungen standen. Die Mischna, die gewöhnlich um 200 n. Chr. mit der Redaktion durch Juda ha-Nasi verbunden wird, stellt einen wichtigen Ausdruck dieser Entwicklung dar; die späteren Talmudim dokumentieren eine umfangreiche Diskussion und Weiterentwicklung rabbinischer Traditionen. Der Übergang vom Judentum des Zweiten Tempels zum rabbinischen Judentum war somit ein längerer geschichtlicher Prozess, in dem ältere Formen religiöser Praxis aufgenommen und unter veränderten Bedingungen neu geordnet wurden.

Auch die Trennung von Judentum und Christentum vollzog sich nicht in einem einzigen historischen Moment (Sanders 1992, 165–175). Die ersten Christusgläubigen waren überwiegend Juden, die Jesus als den von Gott gesandten Messias bekannten; sie nahmen weiterhin an jüdischen Traditionen teil und bezogen sich auf die heiligen Schriften Israels. Mit der zunehmenden Aufnahme nichtjüdischer Anhänger veränderte sich jedoch die Zusammensetzung der Gemeinden, und die Frage nach dem Verhältnis zum jüdischen Gesetz, die Christologie sowie die institutionelle Entwicklung der Gemeinden führten allmählich zu unterscheidbaren religiösen Identitäten. Dieser Prozess verlief regional unterschiedlich und dauerte teilweise bis weit in die Spätantike hinein; deshalb ist auch die Rede von zwei unmittelbar nach 70 n. Chr. getrennten Religionen historisch zu vereinfachend. Das rabbinische Judentum und das Christentum entwickelten sich nicht als zwei bereits vollständig ausgebildete Systeme, die sich plötzlich voneinander lösten, sondern gingen aus einer gemeinsamen antiken Traditionswelt hervor und bildeten in wechselseitiger Abgrenzung unterschiedliche Formen religiösen Lebens aus.

Das Judentum zur Zeit des Zweiten Tempels war damit weder eine bloße Übergangsphase noch eine religiöse Vorstufe des Christentums, sondern eine eigenständige, theologisch produktive und gesellschaftlich vielfältige Epoche, in der ältere Überlieferungen unter veränderten historischen Bedingungen neue Gestalt annahmen. Der Jerusalemer Tempel bildete über Jahrhunderte den Mittelpunkt des religiösen Lebens; zugleich entwickelten sich Formen jüdischer Identität, die über den unmittelbaren Tempelkult hinausreichten. Die Tora, die Auslegung heiliger Schriften, das gemeinschaftliche Gebet und die religiöse Praxis ermöglichten jüdisches Leben in unterschiedlichen Regionen und Kulturen. Die Begegnung mit der hellenistischen Welt führte nicht einfach zu einer Gegenüberstellung von Tradition und Moderne, sondern eröffnete neue Formen religiöser Reflexion, die sich besonders in der Weisheitsliteratur, in der griechischen Bibelübersetzung und im Werk Philos von Alexandria zeigen. Die apokalyptischen und pseudepigraphen Schriften dokumentieren darüber hinaus eine intensive Auseinandersetzung mit der Frage nach Gottes Handeln in der Geschichte; sie entstanden in einer Welt, in der politische Herrschaft, religiöse Identität und die Erfahrung von Unrecht immer wieder neu aufeinandertrafen.

Die Vielfalt von Sadduzäern, Pharisäern, Essenern, apokalyptischen Gemeinschaften und Diasporajuden macht deutlich, dass die jüdische Religion dieser Zeit nicht durch eine einzige theologische Richtung angemessen beschrieben werden kann. Auch die gewaltsame Zerstörung des Garizim-Heiligtums durch Johannes Hyrkanos I. und die in der Forschung diskutierte Spannung zwischen „enochischen“ und „zadokitischen“ Traditionen zeigen, wie sehr religiöse Identität, priesterliche Legitimität und politische Macht ineinandergriffen – ohne dass daraus zwei sauber getrennte „Judentümer“ wurden. Für die theologische Forschung folgt daraus eine methodische Konsequenz: Das frühe Judentum darf weder aus der Perspektive späterer christlicher Dogmatik noch ausschließlich aus der Perspektive späterer rabbinischer Orthodoxie beurteilt werden; seine Schriften, Institutionen und religiösen Gemeinschaften müssen zunächst aus ihren eigenen historischen Voraussetzungen heraus verstanden werden (Sanders 1992, 175–185; Grabbe 2010, 225–235; Collins 2019, 437–459). Erst auf dieser Grundlage lässt sich angemessen beschreiben, welche unterschiedlichen Wege aus der gemeinsamen religiösen Welt des antiken Judentums hervorgingen.
\end{document}